% SHARD-ID: AQM-32-WEYL-HEISENBERG-REPRESENTATIONS-FOR-KT
% SHARD-TITLE: Weyl-Heisenberg Representations for \(k(t)\)
% SHARD-SUMMARY: Defines the algebraic rational-function Heisenberg group using a coefficient character on \(k(t)\).
% SHARD-SUMMARY: Proves the Schrödinger representation on \(\ell^2(k(t))\) is unitary and irreducible.
% SHARD-SUMMARY: Leaves the full LCA, local-field, adelic, and \(k=\mathbb F_3\) examples to the following completion shard.
% SHARD-KEYWORDS: k(t), rational function field, Weyl-Heisenberg, Schrödinger representation, Stone-von Neumann, F3, local field

\section{Weyl-Heisenberg Representations for \texorpdfstring{\(k(t)\)}{k(t)}}
\label{sec:weyl-heisenberg-representations-for-kt}

\subsection{Status and Source Anchors}

This shard uses convention (z).  The LCA uniqueness theorem is sourced from
Prasad, \path{references/heisenberg_weil/prasad_0912_0574_source/SvN.tex},
lines 164--190.  The canonical \(F\times\widehat F\) Weyl system is sourced
from Beny--Crann--Lee--Park--Youn,
\path{references/heisenberg_weil/beny_crann_lee_park_youn_2204_08162_source/main.tex},
lines 461--473.  The locally compact ring theorem and the warning about
countably infinite discrete rings are sourced from Bekka,
\path{references/heisenberg_weil/bekka_2502_00387_source/CCR-GeneralRings-v5.tex},
lines 319--350, 862--881, and 1014.  The local and adelic field background is
registered in \path{references/heisenberg_weil/SOURCES.md} from Poonen and
Kudla.  The algebraic \(k(t)\) calculations below are direct derivations from
the displayed definitions.

\subsection{Lamport Dependency Skeleton}

\[
\begin{array}{rcl}
D1&:&K=k(t),\quad \tau:k\to\mathbb F_p\text{ nonzero},\quad
      \lambda_K(f)=\tau([t^{-1}]f).\\
L1&:&(a,b)\mapsto\chi_K(ab)\text{ is a nondegenerate pairing on }K.\\
D2&:&H_K=\mu_p\times K\times K
      \text{ with cocycle }\chi_K(ba').\\
L2&:&H_K\text{ is a class-two group with center }\mu_p.\\
D3&:&T_a|\xi\rangle=|\xi+a\rangle,\quad
      M_b|\xi\rangle=\chi_K(b\xi)|\xi\rangle.\\
T1&:&\pi(z,a,b)=zT_aM_b\text{ is a unitary irreducible representation
      of }H_K\text{ on }\ell^2(K).\\
T2&:&\text{Stone-von Neumann uniqueness requires the full LCA dual or a
      local/adelic completion.}\\
E1&:&k=\mathbb F_3\text{ gives explicit }\omega\text{-commutation examples.}
\end{array}
\]

\subsection{The Coefficient Character D1}

Let \(k=\mathbb F_q\) have characteristic \(p\), and let \(K=k(t)\).  Fix a
nonzero \(\mathbb F_p\)-linear functional
\[
  \tau:k\longrightarrow\mathbb F_p.
\]
For \(f\in K\), expand \(f\) as a Laurent series at infinity:
\[
  f=\sum_{n\le N} c_n t^n,\qquad c_n\in k.
\]
Only finitely many positive powers occur.  Define
\[
  [t^{-1}]f=c_{-1},
  \qquad
  \lambda_K(f)=\tau([t^{-1}]f).
\]
The associated additive character is
\[
  \chi_K(f)=\exp\!\left(\frac{2\pi i}{p}\lambda_K(f)\right).
\]
The trace \(\operatorname{Tr}_{k/\mathbb F_p}\) is a standard choice for
\(\tau\) when a trace convention is wanted.  The proofs below require only
that \(\tau\ne0\).

\paragraph{Lemma \(L1\).}
The pairing
\[
  B:K\times K\longrightarrow\mathbb C^\times,\qquad
  B(a,b)=\chi_K(ab)
\]
is nondegenerate in each variable.

\emph{Proof.}
Let \(a\in K\) be nonzero.  Since \(\tau\ne0\), choose \(c\in k\) with
\(\tau(c)\ne0\).  Put
\[
  b=\frac{c\,t^{-1}}{a}\in K.
\]
Then \(ab=c\,t^{-1}\), so
\[
  \lambda_K(ab)=\tau(c)\ne0.
\]
Thus \(\chi_K(ab)\ne1\).  This proves that no nonzero \(a\) is killed by all
right pairings.  The same argument with the variables interchanged proves
nondegeneracy in the second variable, because multiplication in \(K\) is
commutative.

\subsection{The Algebraic Heisenberg Group D2--L2}

Let \(\mu_p=\{z\in\mathbb C:z^p=1\}\).  Define
\[
  H_K=\mu_p\times K\times K
\]
with multiplication
\[
  (z,a,b)(z',a',b')
  =
  \bigl(zz'\chi_K(ba'),\,a+a',\,b+b'\bigr).
\]
Here \(a\) is the translation coordinate and \(b\) is the modulation
coordinate.

\paragraph{Associativity.}
Let \(h=(z,a,b)\), \(h'=(z',a',b')\), and \(h''=(z'',a'',b'')\).  The central
factor in \((hh')h''\) is
\[
  zz'z''\chi_K(ba')\chi_K((b+b')a'')
  =
  zz'z''\chi_K(ba'+ba''+b'a'').
\]
The central factor in \(h(h'h'')\) is
\[
  zz'z''\chi_K(b'a'')\chi_K(b(a'+a''))
  =
  zz'z''\chi_K(b'a''+ba'+ba'').
\]
These are equal by additivity of \(\chi_K\).  Therefore the multiplication is
associative.

\paragraph{Identity and inverse.}
The identity is \((1,0,0)\).  A direct multiplication gives
\[
  (z,a,b)^{-1}=\bigl(z^{-1}\chi_K(ba),-a,-b\bigr).
\]

\paragraph{Commutator and center.}
From the multiplication law,
\[
  (1,a,b)(1,a',b')
  =
  \chi_K(ba')(1,a+a',b+b')
\]
and
\[
  (1,a',b')(1,a,b)
  =
  \chi_K(b'a)(1,a+a',b+b').
\]
Hence the commutator is the central element
\[
  [(1,a,b),(1,a',b')]
  =
  \bigl(\chi_K(ba'-b'a),0,0\bigr).
\]
If \((z,a,b)\) is central, then this phase is \(1\) for all \(a',b'\).  Taking
\(b'=0\) gives \(\chi_K(ba')=1\) for all \(a'\), hence \(b=0\) by \(L1\).
Taking \(a'=0\) then gives \(\chi_K(-b'a)=1\) for all \(b'\), hence \(a=0\).
Thus
\[
  Z(H_K)=\mu_p\times\{0\}\times\{0\}.
\]

\subsection{The Schrödinger Representation D3--T1}

Let \(\ell^2(K)\) have orthonormal basis \(\{|\xi\rangle:\xi\in K\}\).  Define
\[
  T_a|\xi\rangle=|\xi+a\rangle,
  \qquad
  M_b|\xi\rangle=\chi_K(b\xi)|\xi\rangle.
\]
Both are unitary: \(T_a\) permutes the basis and \(M_b\) is diagonal with
unit-modulus diagonal entries.  Moreover,
\[
\begin{aligned}
  M_bT_a|\xi\rangle
  &=M_b|\xi+a\rangle                                      \\
  &=\chi_K(b\xi+ba)|\xi+a\rangle                          \\
  &=\chi_K(ba)T_aM_b|\xi\rangle.
\end{aligned}
\]
Therefore
\[
  T_aM_bT_{a'}M_{b'}=\chi_K(ba')T_{a+a'}M_{b+b'}.
\]
The map
\[
  \pi_K:H_K\longrightarrow U(\ell^2(K)),\qquad
  \pi_K(z,a,b)=zT_aM_b
\]
is a unitary representation with central character
\[
  \pi_K(z,0,0)=zI.
\]

\paragraph{Irreducibility.}
Let \(A\) be a bounded operator on \(\ell^2(K)\) commuting with all
\(T_a\) and \(M_b\).  Write
\[
  A_{\xi,\eta}=\langle \xi|A|\eta\rangle.
\]
Commutation with \(M_b\) gives
\[
  \chi_K(b\xi)A_{\xi,\eta}=\chi_K(b\eta)A_{\xi,\eta}
\]
for all \(b,\xi,\eta\).  If \(\xi\ne\eta\), then \(L1\) gives \(b\) such that
\(\chi_K(b(\xi-\eta))\ne1\), so \(A_{\xi,\eta}=0\).  Thus \(A\) is diagonal.
Write \(A_{\xi,\xi}=d_\xi\).  Commutation with \(T_a\) gives
\[
  d_{\xi+a}=d_\xi
\]
for all \(a,\xi\).  Since the additive group of \(K\) acts transitively on
itself by translations, all \(d_\xi\) are equal.  Thus the commutant of
\(\pi_K(H_K)\) is \(\mathbb C I\).  If a closed subspace were invariant, its
orthogonal projection would commute with the representation, hence would be
\(0\) or \(I\).  Therefore \(\pi_K\) is irreducible.

\subsection{Conclusion}

The concrete rational-function answer is:
\[
  H_{\mathbb F_q(t)}=\mu_p\times k(t)\times k(t),
  \qquad
  (z,a,b)(z',a',b')=(zz'\chi_K(ba'),a+a',b+b'),
\]
with irreducible Schrödinger representation on \(\ell^2(k(t))\).  The analytic
Stone-von Neumann answer is separate and is recorded in the following shard.
